%% Font size %%
\documentclass[11pt]{article}

%% Load the custom package
\usepackage{Mathdoc}

%% Numéro de séquence %% Titre de la séquence %%
\renewcommand{\centerhead}{S14- Probabilités - Exercices Divers}

%% Spacing commands %%
\renewcommand{\baselinestretch}{1} \setlength{\parindent}{0pt}

%%\pagestyle{empty}

\begin{document}


\begin{exercice}[2][Événement contraire]

Dans une boîte, il y a 9 boules blanches et 6 boules noires.

On tire une boule au hasard.

\begin{enumerate}
    \item Calculer la probabilité d'obtenir une boule noire.
    \item En déduire la probabilité d'obtenir une boule blanche.
\end{enumerate}

\end{exercice}


\begin{exercice}[2][Compléter avec l'événement contraire]

On lance un dé à 6 faces.

Compléter.

\begin{enumerate}
    \item L'événement contraire de « obtenir un nombre pair » est :
    \dotfill

    \item L'événement contraire de « obtenir un nombre inférieur à 5 » est :
    \dotfill

    \item L'événement contraire de « obtenir un 6 » est :
    \dotfill
\end{enumerate}

\end{exercice}


\begin{exercice}[2][Événement certain ou impossible]

Pour chaque événement, préciser s'il est :

\begin{center}
certain \hspace{1cm} impossible \hspace{1cm} possible
\end{center}

\begin{enumerate}
    \item Obtenir un nombre inférieur à 7 avec un dé à 6 faces.
    \item Obtenir un 9 avec un dé à 6 faces.
    \item Obtenir un nombre pair avec un dé à 6 faces.
    \item Tirer une boule rouge dans un sac qui ne contient que des boules rouges.
\end{enumerate}

\end{exercice}

\end{document}
